\chapter{The Interaction Picture}

\begin{remarkbox}
The previous chapter introduced propagators of free fields. Perturbation theory
uses those free fields as the starting point and treats interactions as a
controlled deformation. The interaction picture is the framework that makes this
separation precise.
\end{remarkbox}

\section{Why We Need a New Picture}

In quantum mechanics one may describe time evolution in several equivalent ways.
The Schroedinger picture puts time dependence in the states. The Heisenberg
picture puts time dependence in the operators. The interaction picture splits
the time dependence between them.

This split is useful when the Hamiltonian has the form
\[
    \hat H = \hat H_0 + \hat H_{\mathrm{int}}.
\]
Here \(\hat H_0\) is the exactly solvable free Hamiltonian, and
\(\hat H_{\mathrm{int}}\) contains interactions. In QFT, \(\hat H_0\) is built
from the free scalar, spinor, or photon fields of the previous chapters. The
interaction part describes processes such as scattering, decay, emission, and
absorption.

The guiding idea is simple:
\[
    \text{solve the free theory exactly, then expand in the interaction}.
\]
The interaction picture is the operator language in which this idea becomes a
systematic perturbation theory.

\section{Schroedinger, Heisenberg, and Interaction Pictures}

Let \(|\Psi_S(t)\rangle\) be a Schroedinger-picture state. It satisfies
\[
    i\frac{\dd}{\dd t}|\Psi_S(t)\rangle
    = \hat H|\Psi_S(t)\rangle.
\]
Operators in this picture are usually time independent unless they have explicit
time dependence.

In the Heisenberg picture, the state is fixed and operators evolve according to
\[
    \hat O_H(t)
    = e^{i\hat H t}\hat O_S e^{-i\hat H t}.
\]
This picture is elegant but difficult for interacting QFT because exact
interacting operators are rarely known.

In the interaction picture, free time evolution is assigned to the operators:
\[
    \hat O_I(t)
    = e^{i\hat H_0 t}\hat O_S e^{-i\hat H_0 t}.
\]
The states evolve only with the interaction. Thus interaction-picture fields are
free fields, even though the states know about the interaction.

\section{Free Fields as the Starting Point}

The scalar interaction-picture field satisfies the free Klein--Gordon equation,
\[
    (\Box+m^2)\hat\phi_I(x)=0.
\]
Its mode expansion is the same free expansion used in Chapter 2:
\[
    \hat\phi_I(x)
    = \int\frac{\dd^3k}{(2\pi)^3}\frac{1}{\sqrt{2\omega_{\vect{k}}}}
      \left(
          \hat a_{\vect{k}}e^{-ik\cdot x}
          +\hat a_{\vect{k}}^\dagger e^{ik\cdot x}
      \right),
\]
where \(k^0=\omega_{\vect{k}}\). Similarly, interaction-picture Dirac and
photon fields are free Dirac and free photon fields.

This is why the propagators of Chapter 6 are the basic building blocks of
perturbation theory. The fields appearing inside the perturbative expansion are
free fields, so their time-ordered products can be evaluated using free-field
correlation functions.

\section{The Interaction Hamiltonian}

The interaction Hamiltonian in the interaction picture is
\[
    \hat H_I(t)
    = e^{i\hat H_0 t}\hat H_{\mathrm{int}}e^{-i\hat H_0 t}.
\]
It is written in terms of interaction-picture fields. For example, a scalar
\(\phi^4\) interaction is described by
\[
    \mathcal{L}_{\mathrm{int}}
    = -\frac{\lambda}{4!}\phi^4.
\]
For this case, the interaction Hamiltonian density is
\[
    \mathcal{H}_{\mathrm{int}}
    = \frac{\lambda}{4!}\phi_I^4,
\]
and
\[
    \hat H_I(t)
    = \int\dd^3x\,\frac{\lambda}{4!}\hat\phi_I(t,\vect{x})^4.
\]
The subscript \(I\) reminds us that the fields evolve freely. The coupling
\(\lambda\) measures the strength of the interaction.

For QED, the interaction Lagrangian is
\[
    \mathcal{L}_{\mathrm{int}}
    = -e\bar\psi\gamma^\mu\psi A_\mu,
\]
so the interaction Hamiltonian is commonly written as
\[
    \hat H_I(t)
    = e\int\dd^3x\,
      \hat{\bar\psi}_I\gamma^\mu\hat\psi_I\hat A_{I\mu}.
\]
This term is responsible for the elementary QED vertex.

\section{Time Evolution Operator}

Interaction-picture states satisfy
\[
    i\frac{\dd}{\dd t}|\Psi_I(t)\rangle
    = \hat H_I(t)|\Psi_I(t)\rangle.
\]
Define the time evolution operator \(\hat U_I(t,t_0)\) by
\[
    |\Psi_I(t)\rangle
    = \hat U_I(t,t_0)|\Psi_I(t_0)\rangle.
\]
Then \(\hat U_I(t,t_0)\) obeys
\[
    i\frac{\partial}{\partial t}\hat U_I(t,t_0)
    = \hat H_I(t)\hat U_I(t,t_0),
    \qquad
    \hat U_I(t_0,t_0)=1.
\]
Formally, this is solved by a time-ordered exponential:
\[
    \hat U_I(t,t_0)
    = T\exp\left(-i\int_{t_0}^{t}\dd t'\,\hat H_I(t')\right).
\]
The time-ordering symbol is essential because \(\hat H_I(t_1)\) and
\(\hat H_I(t_2)\) need not commute at different times.

\section{The Dyson Series}

Expanding the time-ordered exponential gives the Dyson series:
\[
    \hat U_I(t,t_0)
    = 1
      -i\int_{t_0}^{t}\dd t_1\,\hat H_I(t_1)
      +\frac{(-i)^2}{2!}\int_{t_0}^{t}\dd t_1\dd t_2\,
      T\{\hat H_I(t_1)\hat H_I(t_2)\}
      +\cdots.
\]
Equivalently,
\[
    \hat U_I(t,t_0)
    = \sum_{n=0}^{\infty}\frac{(-i)^n}{n!}
      \int_{t_0}^{t}\dd t_1\cdots\dd t_n\,
      T\{\hat H_I(t_1)\cdots\hat H_I(t_n)\}.
\]
Each power of \(\hat H_I\) brings one power of the coupling. In \(\phi^4\)
theory, the \(n\)-th order term contains \(\lambda^n\). In QED, the \(n\)-th
order term contains \(e^n\).

This is the algebraic origin of perturbation theory: physical quantities are
computed as series in coupling constants.

\section{Perturbation Theory and Time-Ordered Correlators}

Suppose we want the interacting time-ordered two-point function of a scalar
field. In the interaction picture it is expressed as
\[
    \langle\Omega|T\hat\phi_H(x)\hat\phi_H(y)|\Omega\rangle
    = \frac{
      \langle0|T\left\{\hat\phi_I(x)\hat\phi_I(y)
      \exp\left[-i\int\dd^4z\,\mathcal{H}_{\mathrm{int}}(z)\right]\right\}|0\rangle
    }{
      \langle0|T\exp\left[-i\int\dd^4z\,\mathcal{H}_{\mathrm{int}}(z)\right]|0\rangle
    }.
\]
Here \(|0\rangle\) is the free vacuum and \(|\Omega\rangle\) is the interacting
vacuum. The denominator removes disconnected vacuum factors. This formula is a
central bridge between operator quantization and Feynman diagrams.

The fields inside the expectation value are free fields. Therefore Wick's
theorem can reduce the time-ordered products to sums of propagators. This will
be the subject of the next chapter.

\section{Vacuum Persistence Amplitude}

The vacuum persistence amplitude is
\[
    \langle0|\hat U_I(\infty,-\infty)|0\rangle.
\]
It is the amplitude for the free vacuum in the far past to become the free
vacuum in the far future in the presence of interactions. Its perturbative
expansion contains vacuum bubble diagrams: diagrams with no external legs.

These vacuum bubbles multiply correlation functions by an overall factor. In
normalized expectation values, they cancel against the denominator. This is why
connected diagrams, rather than all diagrams, carry the main physical content of
scattering amplitudes.

\section{Adiabatic Switching and In/Out States}

Scattering theory compares states in the far past with states in the far future.
The idealization is that interactions are effectively switched off at early and
late times, so that asymptotic states can be described by free particles. This
is expressed schematically by replacing the interaction with
\[
    \hat H_I(t)\longrightarrow e^{-\epsilon |t|}\hat H_I(t),
\]
and taking \(\epsilon\to0^+\) at the end.

This adiabatic switching is a mathematical device. It helps define the
connection between free Fock states and interacting scattering amplitudes. The
states in the far past are called in-states, and those in the far future are
called out-states.

The scattering operator is
\[
    \hat S
    = \hat U_I(\infty,-\infty)
    = T\exp\left(-i\int_{-\infty}^{\infty}\dd t\,\hat H_I(t)\right).
\]
Matrix elements of \(\hat S\) between in- and out-states are the basic objects
of perturbative scattering theory.

\section{Why Exact Interacting QFT Is Difficult}

The free theory is exactly solvable because it is a collection of independent
oscillators. Interactions couple those oscillators. Once interactions are
present, particle number may change, the vacuum may differ from the free vacuum,
and exact operator solutions are usually unavailable.

Perturbation theory avoids solving the interacting theory exactly. It expands
around the free theory and organizes corrections in powers of the coupling.
This strategy is extremely powerful when the coupling is small and the relevant
states are close to free asymptotic particle states.

However, perturbation theory has limits. Strong coupling, bound states,
confinement, and nonperturbative vacuum effects require additional methods. The
interaction picture is therefore not the whole of QFT, but it is the standard
entry point into calculational quantum field theory.

\begin{summarybox}
The interaction picture separates the exactly solvable free Hamiltonian from the
interaction Hamiltonian. Interaction-picture fields are free fields, while
states evolve with \(\hat H_I(t)\). The time evolution operator is a
time-ordered exponential, and its expansion gives the Dyson series. This is the
operator foundation of perturbation theory and the starting point for Wick's
theorem and Feynman diagrams.
\end{summarybox}

\section{Chapter Checklist}

After reading this chapter, one should be able to:
\begin{itemize}
    \item distinguish the Schroedinger, Heisenberg, and interaction pictures;
    \item explain why interaction-picture fields are free fields;
    \item write the interaction Hamiltonian for \(\phi^4\) theory and QED;
    \item define the interaction-picture time evolution operator;
    \item expand the time-ordered exponential as the Dyson series;
    \item explain perturbation theory as an expansion in couplings;
    \item describe the role of the vacuum persistence amplitude;
    \item explain why vacuum bubbles cancel in normalized correlators;
    \item define in-states, out-states, and the scattering operator;
    \item explain why exact interacting QFT is difficult.
\end{itemize}
